\documentclass[11pt]{article}

% ============================================================
% PACKAGES
% ============================================================

\usepackage[T1]{fontenc}
\usepackage{microtype}

\usepackage{graphicx}
\usepackage{caption}
\usepackage{subcaption}
\graphicspath{{documentation/resources/mqa-designs/}{plots/}{plots_prefill_corrected/}{plots_reuse/}{images/}{./}}

\usepackage[style=apa]{biblatex}
\addbibresource{bib.bib}

\usepackage{geometry}
\usepackage{amsmath}
\usepackage{amssymb}
\usepackage{booktabs}
\usepackage{enumitem}
\usepackage{float}
\usepackage[hidelinks]{hyperref}

% ============================================================
% FORMATTING
% ============================================================

\geometry{
    a4paper,
    total={170mm,257mm},
    left=20mm,
    top=20mm,
}

\setlength{\parindent}{0pt}
\setlength{\parskip}{0.45em}

\linespread{1.03}

% ============================================================
% TITLE
% ============================================================

\title{
\textbf{
Performance Modeling of a Key-Value-Stationary
Systolic Array for Multi-Query Attention Inference
}
\\[1ex]
\large Final Report
}

\author{
Hudson Strauss
}

\date{\today}

% ============================================================
% DOCUMENT
% ============================================================

\begin{document}

\maketitle

\begin{center}
High-Performance Computer Architecture and Systems \\
Santa Clara University
\end{center}

\vspace{0.5em}

% ============================================================
\section{Introduction}

Transformer inference has made memory-efficient attention a central
systems problem. During autoregressive decoding, each generated token
repeatedly accesses the key-value (KV) cache, causing memory traffic
to scale with sequence length. As context windows increase, inference
becomes increasingly memory-bound rather than compute-bound.

Multi-Query Attention (MQA) reduces this overhead by sharing keys and
values across multiple query heads, substantially lowering KV-cache
footprint compared to standard multi-head attention
\parencite{ainslie2023gqa}. However, most accelerators still map
attention onto generic matrix multiplication hardware despite the
streaming nature of autoregressive decoding.

Traditional systolic-array dataflows --- output-stationary,
weight-stationary, and input-stationary --- are optimized for
dense GEMM workloads and do not exploit the fact that during
MQA inference, queries stream continuously while keys and
values remain relatively fixed across heads.

This report presents an analytical performance model of a
KV-stationary systolic-array mapping for MQA inference,
validated against SCALE-Sim cycle-accurate simulation.
We sweep sequence length, merge-extension depth, and
interleaved-MAC count, comparing the architecture against
both a roofline baseline and a FlashAttention-style fused
reference on equal silicon.
The key finding, under a SCALE-Sim-validated wavefront pipeline-fill
model, is that prefill on this architecture is \emph{DRAM-bound}: it
matches a fused FlashAttention baseline on cycles at equal silicon
(\(\sim\!1\times\)), and the lower-MAC count must be sized to memory
bandwidth (\(lmc \approx \text{BW}/128\)) rather than maximized ---
over-provisioning MACs simply leaves them idle against the memory wall.
An earlier dwell-fill model had suggested a \(\sim\)1.75$\times$
per-area advantage; SCALE-Sim measurement of the pipeline fill
(\(\sim\)2.5 vs an assumed 135 cycles/column) overturns that and yields
the bandwidth-limited picture above.
Autoregressive decode is not validated in this work;
the analytical model projects significant DRAM savings
from holding the KV cache on-chip, and a two-pass decode
extension is identified as the primary path to reducing
chip area for practical deployment.

% ============================================================
\section{Problem Statement}

The goal of this project is to evaluate a KV-stationary systolic-array
dataflow for MQA inference and determine the conditions under which
it outperforms conventional systolic mappings.

For each incoming query vector \(Q\), attention decoding proceeds
in three stages:

\vspace{-0.5em}

\[
s_i = Q \cdot K_i
\]

\[
a_i =
\frac{e^{s_i}}
{\sum_j e^{s_j}}
\]

\[
O = \sum_i a_i V_i
\]

\vspace{-0.5em}

where:

\begin{center}
\begin{tabular}{ll}
\(Q\) & Incoming query vector \\
\(K_i\) & Key stored in PE column \(i\) \\
\(V_i\) & Value vector stored in PE column \(i\) \\
\(s_i\) & Local attention score \\
\(a_i\) & Normalized attention weight \\
\end{tabular}
\end{center}

A critical bottleneck in conventional hardware is the
softmax normalization step: computing
\(\sum_j e^{s_j}\) requires a complete pass over all scores,
forcing either (a) the full \(T \times T\) score matrix to be
materialized in DRAM before normalization, or (b) a two-pass
on-chip tiling strategy as in FlashAttention.
At \(T = 8{,}192\) with \(H = 64\) heads, materializing scores
requires writing and re-reading \textbf{16\,GB} of intermediate
data per prefill pass.

The proposed architecture resolves this by computing the
online softmax update inline at each PE column, eliminating
the score matrix entirely.

% ============================================================
\section{Proposed Architecture}

\subsection{Design Evolution}

The architecture went through four distinct iterations before
arriving at the final configuration.
Each iteration was motivated by a measurable inefficiency in
the previous design.

\textbf{Iteration 1 --- Basic 1D KV-stationary pipeline.}
The starting point is a single row of processing elements,
each storing one \((K_i, V_i)\) pair.
A query vector enters at column~0 and streams left-to-right,
computing a running attention score at each column using the
online softmax update described in \S2.
This eliminates score-matrix DRAM traffic but immediately
exposes a pipeline drain problem: after the query packet enters
the last column it must still traverse all remaining
\(\text{cols} - 1\) columns before the result is available.
At \(T = 8{,}192\) with no partitioning, drain accounts for
essentially 100\% of total latency.
A second problem is also present but less obvious at this stage:
the dot product in each column takes \(d = 128\) cycles,
while the upper Running Attention PE needs only 7 cycles to
process the resulting score.
The upper PE therefore sits idle for 121 of every 135 column
cycles --- a utilization of \(\approx\)5\%.

\textbf{Iteration 2 --- 2D multi-head extension.}
The natural extension is to assign each of the \(H = 64\)
query heads its own row lane.
MQA's defining property --- all heads share a single K and V
cache --- means the same column data is broadcast vertically
to all rows simultaneously.
KV is loaded once from DRAM regardless of head count.
This fills the Y-axis of the array and improves total throughput
\(H\times\), but does not address the drain problem or the
5\% upper PE utilization.

\textbf{Iteration 3 --- Merge extensions to shorten the column traversal.}
A query's latency is dominated by the number of columns it must
ripple through. Partitioning the KV cache across \(2^n\) sub-arrays
stacked in the \emph{row} direction lets each sub-array hold
\(T / 2^n\) tokens and run in parallel, so a query traverses only
\(T / 2^n\) columns instead of \(T\) --- trading the vertical dimension
for the horizontal one.
A binary tree of inline merge PEs (Figure~\ref{fig:merge_reduction})
recombines the partial \((m, l, O)\) softmax states with no
post-pipeline overhead.
With \(n = 3\) the traversal drops \(8\times\) for \(8\times\) more PE
rows; the PE \emph{count} is unchanged, as the array is merely reshaped
taller and narrower.
Whether this pays off is regime-dependent: under the SCALE-Sim-validated
wavefront model (\S\ref{sec:pipeline}) the column traversal is already
cheap and prefill is DRAM-bound, so merge extensions help only when the
design is compute-bound --- the quantitative comparison is deferred to
\S\ref{sec:equalsilicon}.

\textbf{Iteration 4 --- Interleaved MACs (\(lmc\)) to improve PE utilization.}
Profiling the cycle breakdown revealed that merge extensions
were treating the symptom (drain) rather than the root cause
(idle upper PEs).
With a single MAC per column the upper PE fires once every
135 cycles but needs only 7 --- a 95\% idle rate.
The fix is to time-interleave \(lmc\) independent MAC units
within each column, each servicing a different query packet
offset by \(\lceil d / lmc \rceil\) cycles.
With \(lmc = 16\): a new query packet starts every 8 cycles
instead of every 128, and upper PE utilization rises from
5\% to 87.5\%.
Crucially, this requires \textit{no extra PE rows} ---
the interleaving is entirely within each existing column.
The quantitative payoff of \(lmc\) (and how it interacts with \(n\))
depends on memory bandwidth: the corrected, SCALE-Sim-validated
comparison is given in \S\ref{sec:equalsilicon}. At commodity
bandwidth prefill is DRAM-bound and \(lmc\) must be \emph{sized} to it
(\(lmc \approx \text{BW}/128\)) rather than maximized; with abundant
bandwidth the array is compute-bound and larger \(lmc\) pays off.

\textbf{Design conclusion.}
Merge extensions buy a shorter column traversal at the cost of rows;
interleaved MACs buy utilization at no row cost. Which combination is
best is \emph{bandwidth-dependent} (\S\ref{sec:equalsilicon}): at
commodity bandwidth prefill is DRAM-bound and both knobs largely wash
out (size \(lmc\) to bandwidth; \(n = 0\) suffices), whereas with
abundant bandwidth the design is compute-bound and both knobs ---
including merge extensions --- contribute. An earlier dwell-fill timing
model had instead singled out \(lmc = 16, n = 0\) as universally
Pareto-optimal; \S\ref{sec:pipeline} explains why that conclusion does
not survive the corrected fill model.

\subsection{Array Layout and Dataflow}

The array is organized as a 2D grid with
\textit{rows} = parallel query lanes and
\textit{columns} = KV-stationary token positions.
For a configuration with \(H\) heads and merge-extension
depth \(n\), the physical dimensions are:

\[
\text{array\_rows} = H \cdot 2^n, \qquad
\text{array\_cols} = \left\lfloor T \,/\, 2^n \right\rfloor
\]

Query vectors enter at column 0 and stream left-to-right.
Keys and values are loaded once into the columns from DRAM
at the start of each sequence and remain resident throughout
the pass.

Each column contains two pipeline stages, illustrated in
Figure~\ref{fig:kv_stationary_architecture}:

\textbf{Stage 1 --- Lower MAC array.}
\(lmc\) interleaved serial MAC units compute
\(s_i = Q \cdot K_i\) over \(d\) cycles.
With \(lmc > 1\), multiple query packets share the column
in time-interleave, each offset by
\(\lceil d / lmc \rceil\) cycles (Figure~\ref{fig:pe_utilization}).

\textbf{Stage 2 --- Upper Running Attention PE.}
Receives the score \(s_i\) from Stage~1 and the running
state \((m, l, O)\) from the left neighbour, then executes
the online softmax update:

\[
m_{\text{out}} = \max(m_{\text{in}},\, s_i)
\]
\[
l_{\text{out}} = l_{\text{in}}\,e^{(m_{\text{in}} - m_{\text{out}})}
               + e^{(s_i - m_{\text{out}})}
\]
\[
O_{\text{out}} = O_{\text{in}}\,e^{(m_{\text{in}} - m_{\text{out}})}
               + V_i\,e^{(s_i - m_{\text{out}})}
\]

After the query packet has traversed all columns,
the final output is normalized:
\(\text{Output} = O_N / l_N\).

\subsection{Pipeline Timing Model}
\label{sec:pipeline}

Let \(d = 128\) (head dimension), \(w = 128\) (upper PE MAC width,
fully parallel), and \(\tau_{\exp} = 4\) cycles (exp lookup). The two
stages have the following latencies:
\[
\text{upper\_pe} = \tau_{\exp} + \left\lceil \tfrac{3d}{w} \right\rceil = 7
\;\text{cycles},
\qquad
\text{column\_dwell} = d + \text{upper\_pe} = 135 \;\text{cycles}.
\]
Here \(\text{column\_dwell}\) is the latency for a \emph{single} packet to
finish at one column (its full \(d\)-cycle dot product plus the upper PE).
We stress that these per-stage cycle costs are \emph{modeling assumptions}, not
figures from a hardware implementation: \(\tau_{\exp} = 4\) (and hence
\(\text{upper\_pe} = 7\)) is an assumed exp-lookup latency, and the
\(\lceil 3d/w \rceil\) accumulation term assumes a fully parallel \(w = d\)-wide
upper PE. A real LUT, exp datapath, or narrower MAC width would change these
constants. SCALE-Sim validates the \(Q\!\cdot\!K\) traversal structure
(\S\ref{sec:wavefront}); it does \emph{not} validate the upper-PE latencies,
which remain analytical throughout (see \S\ref{sec:gaps}).
The lower MAC interleaves \(lmc\) query packets, emitting one score every
\(\text{packet\_stagger} = \lceil d / lmc \rceil\) cycles; a shallow score
buffer then paces delivery to the upper PE at
\[
\text{effective\_stagger} = \max(\text{packet\_stagger},\,\text{upper\_pe}),
\]
so the upper PE receives exactly one score per \(\text{effective\_stagger}\)
cycles. Its utilization, \(\text{upper\_pe} / \text{effective\_stagger}\), rises
from \(\approx\!5\%\) at \(lmc = 1\) (stagger \(= 128\)) to \(87.5\%\) at
\(lmc = 16\) (stagger \(= 8\)).

The key timing property is that every column owns a \emph{dedicated} lower MAC
and a stationary \(K_i\). For a single query packet the \(C\) dot products
\(s_0 \dots s_{C-1}\) are therefore independent: they compute as a pipelined
\emph{wavefront} as \(Q\) ripples right, one column per register hop. The
\(d\)-cycle reduction latency is paid \emph{once} as pipeline fill --- not at
every column --- after which the binding rate is the online-softmax state chain
running at \(\text{effective\_stagger}\). For a tile with \(R\) active rows,
\(C\) active columns, and \(P\) packets per row:
\[
\text{tile\_cycles}
= \underbrace{d}_{\text{fill}}
+ \underbrace{(R + C - 1)\cdot\text{effective\_stagger}}_{\text{traversal}}
+ \underbrace{(P - 1)\cdot\text{effective\_stagger}}_{\text{packet stream}}.
\]

A naive alternative --- assuming a packet \emph{dwells} the full
\(\text{column\_dwell}\) at each column before advancing, giving
\((R + C - 1)\cdot\text{column\_dwell}\) --- would overcount the pipeline fill by
\(\sim\!50\times\). We reject it on measurement: SCALE-Sim reports the per-column
fill of the \(Q\!\cdot\!K\) stage at \(\sim\!2.5\) cycles, not \(135\)
(\S\ref{sec:wavefront}), confirming the wavefront model. A direct consequence,
developed in \S\ref{sec:wavefront}, is that prefill on this array is
\emph{DRAM-bound}; all cycle counts reported in this paper use the wavefront fill.

\subsection{Merge Extensions}

A query packet's latency is dominated by the \emph{horizontal} traversal: it must
ripple through all \(C\) columns of its sub-array, so the pipeline depth
\(R + C - 1\) (\S\ref{sec:pipeline}) grows with the column count for long
sequences. Merge extensions cut this total cycle time by trading the vertical
dimension for the horizontal one. The KV cache is split across \(2^n\)
sub-arrays stacked in the \emph{row} direction
(Figure~\ref{fig:merge_reduction}): each sub-array holds \(T/2^n\) tokens and is
processed \emph{in parallel} on its own rows, so a query traverses only
\(T/2^n\) columns instead of \(T\). A binary tree of inline merge PEs then
recombines the \(2^n\) partial \((m, l, O)\) states in \(n\) extra pipeline
levels. The net effect is fewer horizontal layers to pass through --- the
depth-dominating column term shrinks by \(2^n\) while the row count grows by
\(2^n\) --- lowering total latency, paid for with \(2^n\times\) more PE rows
running concurrently.

We introduce merge extensions here for completeness, as they were a central
design lever under the original (dwell-fill) timing assumption. Under the
SCALE-Sim-validated wavefront model, however, the column traversal is already
cheap and prefill is DRAM-bound (\S\ref{sec:pipeline}, \S\ref{sec:wavefront}),
so merge extensions yield little for prefill --- their benefit is confined to
the compute-bound regime (low \(lmc\) or high bandwidth), as quantified in
\S\ref{sec:wavefront}.

% ============================================================
% ARCHITECTURE FIGURES
% ============================================================

\begin{figure}[H]
    \centering
    % Row 1
    \begin{subfigure}[t]{0.48\textwidth}
        \centering
        \includegraphics[width=\textwidth]{AttentionPE.png}
        \caption{\textbf{Running Attention PE.}
        Query \(Q_0\) (blue) streams left-to-right past stationary
        \(K_i, V_i\) (red). Stage~1 computes \(s_i = Q \cdot K_i\);
        Stage~2 updates the online softmax state \((m, l, O)\)
        and forwards it to the next column.}
        \label{fig:attention_pe}
    \end{subfigure}
    \hfill
    \begin{subfigure}[t]{0.48\textwidth}
        \centering
        \includegraphics[width=\textwidth]{FinalOutVector.png}
        \caption{\textbf{Final output vector.}
        After traversing all columns the accumulated running state
        is normalised: \(\text{Output} = O_N / l_N\).
        No score matrix is written to DRAM at any point.}
        \label{fig:final_output}
    \end{subfigure}

    \vspace{0.8em}

    % Row 2
    \begin{subfigure}[t]{0.48\textwidth}
        \centering
        \includegraphics[width=\textwidth]{MultiHeaded.png}
        \caption{\textbf{Multi-headed MQA mapping.}
        Each query head occupies a dedicated row lane.
        All heads share the same stationary \(K_i, V_i\) columns
        (MQA property), so KV is loaded once regardless of head count.}
        \label{fig:multiheaded}
    \end{subfigure}
    \hfill
    \begin{subfigure}[t]{0.48\textwidth}
        \centering
        \includegraphics[width=\textwidth]{SystolicSingle.png}
        \caption{\textbf{Single-head systolic operation.}
        The query vector streams through a MAC row, accumulating
        \(Q \cdot K\) element-by-element. The completed score
        propagates upward to its Running Attention PE column.}
        \label{fig:systolic_single}
    \end{subfigure}

    \caption{KV-stationary MQA dataflow. Keys and values remain
    stationary in array columns; query vectors stream left-to-right.
    Online softmax is computed inline at every column, eliminating
    the \(H \times T^2\) score-matrix DRAM traffic of an unfused baseline.}
    \label{fig:kv_stationary_architecture}
\end{figure}

\begin{figure}[H]
    \centering
    \includegraphics[width=0.82\textwidth]{MergeReductionGraph.png}
    \caption{\textbf{High-\(T\) merge reduction graph} (\(n = 2\) levels shown).
    Four sub-arrays (red KV rows) each process \(T/4\) tokens independently.
    A two-level binary tree of inline merge PEs (blue) combines partial
    \((m, l, O)\) states, cutting per-sub-array drain from \(T-1\) to
    \(T/2^n - 1\) steps at the cost of \(2^n\times\) more PE rows.}
    \label{fig:merge_reduction}
\end{figure}

% ============================================================
\section{Experimental Methodology}
\label{sec:methodology}

\subsection{Hybrid Analytical and SCALE-Sim Model}

The performance model combines an analytical pipeline timing
model with SCALE-Sim cycle-accurate validation for the
compute-dominant GEMM components.

\textbf{SCALE-Sim validated components:}
\begin{itemize}[noitemsep]
    \item Lower MAC unit --- the \(Q \cdot K\) dot-product GEMM,
          mapped as shape \((H \cdot Q_T / lmc) \times d \times \text{eff\_cols}\)
          on a weight-stationary \(H \times \text{eff\_cols}\) array.
    \item FlashAttention QK tile and AV tile --- both mapped as
          independent GEMMs on a \(64 \times 64\) reference array.
\end{itemize}

\textbf{Analytical components:}
\begin{itemize}[noitemsep]
    \item Upper PE latency: \(\tau_{\exp} + \lceil 3d / w \rceil = 7\) cycles.
    \item Full pipeline timing: fill, traversal, and packet-stream
          terms computed from the wavefront model (\S\ref{sec:pipeline}).
    \item Merge-tree synchronization overhead.
\end{itemize}

The corrected pipeline total substitutes the SCALE-Sim lower-MAC
measurement for the analytical estimate:

\[
\text{corrected\_cycles}
= \text{pipeline\_cycles} - \text{analytical\_mac} + \text{ss\_mac\_cycles}
\]

\subsection{Validation Error Analysis}

The analytical lower-MAC bound is:

\[
\text{analytical\_mac} = \left\lceil Q_T \cdot d \;/\; lmc \right\rceil
\]

This counts pure computation cycles and excludes pipeline drain ---
the time the last packet takes to flush through the remaining
\(\text{eff\_cols} - 1\) columns after the array stops accepting
new work.
SCALE-Sim includes this drain, so its cycle count is always
\(\geq\) the analytical prediction.
The percentage error is:

\[
\delta\% = 100 \cdot
\frac{\text{ss\_mac\_cycles} - \text{analytical\_mac}}
     {\text{analytical\_mac}}
\]

As sequence length grows, the computation term dominates and
drain becomes a smaller fraction of total cycles;
error converges toward zero at large \(T\).
Observed errors for the prefill configurations:
\(lmc = 1\): \(<3\%\) at \(T \geq 512\);
\(lmc = 16\): \(4\)--\(9\%\) at \(T \geq 1{,}024\).

We emphasize that this \(\delta\%\) is not modeling noise --- it \emph{is} the
physical pipeline fill, and measuring it settles the dwell-vs-wavefront question
of \S\ref{sec:pipeline}. Across \(40\) SCALE-Sim \(Q\!\cdot\!K\) GEMMs spanning
\(lmc \in \{1,2,4,8,16\}\), the drain corresponds to only \(\sim\!2.5\) cycles per
column (and is independent of \(lmc\), being set by the systolic geometry). The
dwell model's \(135\) cycles per column is therefore wrong by \(\sim\!50\times\);
the wavefront fill is the validated one. We develop this and its consequences in
\S\ref{sec:wavefront}.

Decode configurations show errors of hundreds to thousands
of percent because the analytical lower-MAC formula counts only
\(d = 128\) MAC cycles while SCALE-Sim measures the full
\(\text{eff\_cols}\)-column pipeline traversal.
Decode behavior is therefore \textbf{not validated} in this work.
Validating the full decode pipeline --- including the interaction
between on-chip KV retention, SRAM capacity constraints, and
multi-pass scheduling --- is left as future work (\S\ref{sec:future}).

\subsection{Fixed Hardware Parameters}

All experiments use the following hardware constants,
chosen to model a realistic edge-inference accelerator:

\begin{center}
\begin{tabular}{ll}
\toprule
Parameter & Value \\
\midrule
Query heads \(H\) & 64 \\
Head dimension \(d\) & 128 \\
Precision & FP16 (2 bytes/element) \\
DRAM bandwidth & 512 B/cycle \\
Upper PE MAC width \(w\) & 128 (fully parallel) \\
Exp lookup latency \(\tau_{\exp}\) & 4 cycles \\
\bottomrule
\end{tabular}
\end{center}

\subsection{Parameter Sweep}

A total of 29 configurations were swept, covering all
combinations of:

\begin{center}
\begin{tabular}{ll}
\toprule
Parameter & Values \\
\midrule
Sequence length \(T\) & 128, 512, 1024, 2048, 4096, 8192 \\
Merge extensions \(n\) & 0 (flat array), 3 (binary tree) \\
MACs per PE column (\(lmc\)) & 1, 16 \\
Mode & Prefill (\(Q_T = T\)), Decode (\(Q_T = 1\)) \\
\bottomrule
\end{tabular}
\end{center}

SCALE-Sim validation was run where memory constraints
permitted (operand matrices below 500\,MB).
Configurations exceeding this limit are reported from the
analytical model only and are flagged accordingly.

\subsection{Experiments Conducted}

Three distinct experiments were run, each targeting a different
layer of the model.

\textbf{Experiment 1 --- Lower MAC validation (29 configurations).}
Each configuration maps the Q·K dot-product GEMM onto SCALE-Sim
as a weight-stationary array of shape \(H \times \text{eff\_cols}\)
with GEMM dimensions \(M = H \cdot Q_T / lmc\), \(N = \text{eff\_cols}\),
\(K = d\).
The analytical prediction \(\lceil Q_T \cdot d / lmc \rceil\) (pure compute)
is compared against SCALE-Sim compute cycles;
the delta \(\delta\%\) is the \emph{pipeline fill} SCALE-Sim measures beyond
that bound --- the quantity that decides the dwell-vs-wavefront question of
\S\ref{sec:pipeline} (analyzed in \S\ref{sec:wavefront}), not a modeling error.
Of 29 configurations, 16 prefill runs completed successfully;
decode runs were excluded after all 12 showed divergence exceeding
200\% (see \S4.2).
Selected prefill validation results are shown in
Table~\ref{tab:mac_validation}.

\begin{table}[H]
\centering
\caption{Lower MAC SCALE-Sim validation --- selected prefill configurations.
\(\delta\%\) = SCALE-Sim pipeline fill as a percentage of the analytical MAC
bound (the measured wavefront fill, \S\ref{sec:wavefront}).
Mac frac = SCALE-Sim MAC cycles as a fraction of the (dwell-model) total
pipeline cycles.}
\label{tab:mac_validation}
\begin{tabular}{rrrrrrr}
\toprule
\(T\) & \(n\) & \(lmc\) & Analytical & SCALE-Sim & \(\delta\%\) & Mac frac \\
\midrule
128  & 0 &  1 & 16,384 & 16,891 & 3.1\%  & 40.2\% \\
128  & 0 & 16 &  1,024 &  1,531 & 49.5\% &  5.7\% \\
512  & 0 &  1 & 65,536 & 66,811 & 1.9\%  & 46.7\% \\
512  & 0 & 16 &  4,096 &  5,371 & 31.1\% &  6.6\% \\
1024 & 0 & 16 &  8,192 & 10,491 & 28.1\% &  6.8\% \\
2048 & 0 & 16 & 16,384 & 20,731 & 26.5\% &  6.9\% \\
2048 & 3 & 16 & 16,384 & 17,147 &  4.7\% &  3.6\% \\
4096 & 3 & 16 & 32,768 & 34,043 &  3.9\% &  3.9\% \\
8192 & 3 & 16 & 65,536 & 67,835 &  3.5\% &  4.0\% \\
\bottomrule
\end{tabular}
\end{table}

Two patterns are visible, and both follow from the fill being small and
roughly constant per column (\(\sim\!2.5\) cycles, \S\ref{sec:wavefront})
rather than from any error that ``converges.''
First, for \(lmc = 1\) the compute bound is large (\(128\) cycles per packet),
so the fixed fill is a small fraction --- \(\delta\% < 3\%\) at \(T \geq 512\).
Second, for \(lmc = 16\) the compute bound is \(8\times\) smaller, so the same
small fill is a larger \emph{fraction} at small \(T\) (\(49\%\) at \(T = 128\));
as \(T\) grows the throughput term grows and the fraction falls below \(7\%\)
by \(T = 1{,}024\).
The Mac frac column reports the lower MAC as only 4--7\% of the
\emph{dwell-model} pipeline total; because that total over-counts the fill by
\(\sim\!50\times\) (\S\ref{sec:pipeline}), it should \emph{not} be read as
evidence that the upper PE or pipeline fill is the bottleneck. Under the
corrected wavefront model the fill is cheap and prefill is DRAM-bound
(\S\ref{sec:wavefront}).

\textbf{Experiment 2 --- Prefill vs FlashAttention sweep (24 configurations).}
KV-stat is compared against a FlashAttention-style fused reference
across all \((T, n, lmc)\) combinations.
Both architectures use a \(64 \times 64\) FlashAttention tile
(SRAM-constrained to \(H = 64\) tile size).
The KV-stat corrected cycle count from Experiment~1 is used where
SCALE-Sim was available; the analytical pipeline total is used
otherwise.
Results are reported as raw speedup, MAC ratio, and per-physical-MAC
speedup, for both the DRAM-bound and compute-only bandwidth regimes
(Table~\ref{tab:prefill_results}).
Of 24 configurations, 15 use SCALE-Sim-corrected KV-stat cycles;
9 use analytical estimates only (flagged in the table).

\textbf{Experiment 3 --- Cycle breakdown decomposition.}
For each validated configuration the total pipeline cycles are
decomposed into three additive terms:
(i) SCALE-Sim lower MAC cycles (\texttt{ss\_lower\_mac}),
(ii) within-row stagger cost \((P-1) \cdot \text{packet\_stagger}\),
and (iii) upper PE and pipeline fill/drain overhead.
This decomposition is \emph{compute-only} (SCALE-Sim runs at infinite
bandwidth): for \(n=3, lmc=16\) the lower-MAC and stagger terms are
roughly balanced, while for \(n=0, lmc=16\) at large \(T\) the within-row
stagger term dominates the compute. Once DRAM bandwidth is modelled,
however, this entire compute total falls below the memory floor and
prefill is DRAM-bound (\S\ref{sec:wavefront}) --- so the throughput
bottleneck is DRAM, not the upper PE.

% ============================================================
% PE UTILIZATION FIGURE
% ============================================================

\begin{figure}[H]
    \centering
    \includegraphics[width=0.82\textwidth]{PEUtilization.png}
    \caption{\textbf{PE utilization with \(lmc = 16\) interleaved MACs.}
    Each column is fed by two streams: a \emph{Q stream} of query packets entering
    from the left, and a \emph{K stream} supplying the stationary key \(K_i\) at
    the top of each column. The sixteen MAC units time-multiplex sixteen query
    packets (\(Q_0\)--\(Q_{15}\), with \(Q_{16}\)--\(Q_{31}\) queued behind), each
    staggered by \(\lceil d/16 \rceil = 8\) cycles, and every MAC emits a score
    \(Q\cdot K\) up to its Running Attention PE. All sixteen share the column's
    single on-chip K register, so no \(K\) duplication is needed. The initiation
    interval drops from 128 to 8 cycles, raising upper-PE utilisation from
    \(\approx\)5\% to \(\approx\)87.5\%.}
    \label{fig:pe_utilization}
\end{figure}

% ============================================================
\section{Results}
\label{sec:results}

\subsection{Baseline Characterization}
\label{subsec:baseline}

Three baseline variants are used throughout the comparison,
each representing a distinct hardware or software assumption.

\textbf{Variant~1 --- Original \(64 \times 64\), unfused.}
A 4,096 MAC weight-stationary array running two sequential GEMMs
(Q·K\(^\top\), then scores·V) with the full score matrix written
to and read back from DRAM between them.
This is the default SCALE-Sim reference configuration.

\textbf{Variant~2 --- MAC-normalized, unfused.}
Same two-GEMM structure with the array resized to match
KV-stat's MAC count: \(8{,}192 \times (T/8)\) = 8.4M MACs
at \(T = 8{,}192\).
This isolates architectural efficiency from raw compute budget.

\textbf{Variant~3 --- MAC-normalized, fused.}
Same large array as Variant~2 with FlashAttention-style tiling:
no score matrix written to DRAM.
This is the fairest possible comparison --- equal MACs,
equal memory discipline.

\subsubsection*{The \(64 \times 64\) Baseline Is Compute-Bound at Every \(T\)}

The original \(64 \times 64\) array is perpetually compute-bound:
its 4,096 MACs cannot keep pace with the prefill workload at any
tested sequence length.

\begin{table}[H]
\centering
\caption{Original \(64 \times 64\) unfused baseline roofline analysis
across all sequence lengths (prefill, \(H=64\), \(d=128\), FP16).}
\label{tab:baseline_orig}
\begin{tabular}{rrrrr}
\toprule
\(T\) & Compute cycles & Memory cycles & Estimated cycles & Bottleneck \\
\midrule
128   &      65,536 &      16,512 &      65,536 & COMPUTE \\
512   &   1,048,576 &     164,352 &   1,048,576 & COMPUTE \\
1,024 &   4,194,304 &     590,848 &   4,194,304 & COMPUTE \\
2,048 &  16,777,216 &   2,230,272 &  16,777,216 & COMPUTE \\
4,096 &  67,108,864 &   8,654,848 &  67,108,864 & COMPUTE \\
8,192 & 268,435,456 &  34,086,912 & 268,435,456 & COMPUTE \\
\bottomrule
\end{tabular}
\end{table}

At \(T = 8{,}192\) the compute bottleneck is \(7.8\times\) worse
than the memory bottleneck.
Comparing KV-stat (memory-bound, 8.4M MACs) against this baseline
guarantees a large speedup regardless of architectural quality;
the 511$\times$ speedup figure is almost entirely a MAC-count artifact.

\subsubsection*{Score-Matrix Dominance}

Once MACs are equalized, the unfused baseline immediately flips
to memory-bound --- and the score matrix accounts for essentially
all of the DRAM traffic.
At \(T = 8{,}192\):

\begin{table}[H]
\centering
\caption{DRAM traffic breakdown for the MAC-normalized unfused baseline
at \(T = 8{,}192\), \(H=64\), \(d=128\), FP16.}
\label{tab:dram_breakdown}
\begin{tabular}{lrr}
\toprule
Traffic component & Size (MB) & \% of total \\
\midrule
Score matrix (write + read) & 17,179.9 & 98.4\% \\
Q reads                     &    134.2 &  0.77\% \\
Output writes               &    134.2 &  0.77\% \\
K reads                     &      2.1 &  0.01\% \\
V reads                     &      2.1 &  0.01\% \\
\midrule
\textbf{Total}              & \textbf{17,452.5} & \\
\bottomrule
\end{tabular}
\end{table}

The score matrix scales as
\(H \times T^2 \times \text{bpe} \times 2 = 64 \times 8192^2 \times 2 \times 2\)
bytes.
This is why the unfused MAC-normalized baseline is \(65\times\)
slower than KV-stat: 17.45\,GB vs 0.27\,GB at the same bandwidth
(Figure~\ref{fig:dram_breakdown}).

\begin{figure}[H]
    \centering
    \includegraphics[width=0.72\textwidth]{fig18_dram_breakdown.png}
    \caption{DRAM traffic at \(T = 8{,}192\). The unfused baseline is dominated by
    the \(H\times T^2\) score matrix (98.4\% of \(17.45\,\)GB); both fused designs
    --- KV-stat and FlashAttention --- eliminate it, dropping to \(0.27\,\)GB. The
    \(65\times\) advantage is score-matrix elimination, not the KV-stationary
    dataflow itself.}
    \label{fig:dram_breakdown}
\end{figure}

\subsubsection*{Fused Baseline Closes the Gap}

The fused (FlashAttention-style) baseline eliminates score-matrix
traffic entirely, dropping total DRAM to 272.6\,MB --- identical
to KV-stat.
Arithmetic intensity jumps from 63 to 4,033 MACs/byte.
Table~\ref{tab:ai_comparison} compares all four architectures
at \(T = 8{,}192\).

\begin{table}[H]
\centering
\caption{Arithmetic intensity at \(T = 8{,}192\), \(H=64\), \(d=128\), FP16.
KV-stat has \(2\times\) the MACs of the baselines (dot product + value accumulation)
but the same 0.27\,GB DRAM as the fused baseline, giving \(2\times\) higher AI.}
\label{tab:ai_comparison}
\begin{tabular}{lrrr}
\toprule
Configuration & Total MACs & Total DRAM & AI (MACs/byte) \\
\midrule
\(64 \times 64\) unfused    & 1.10T & 17.45\,GB & 63 \\
MAC-norm unfused            & 1.10T & 17.45\,GB & 63 \\
MAC-norm fused              & 1.10T &  0.27\,GB & 4,033 \\
KV-stat \(n=3,\;lmc=16\)   & 2.20T &  0.27\,GB & 8,066 \\
\bottomrule
\end{tabular}
\end{table}

\subsubsection*{Three-Way Speedup Summary}

Table~\ref{tab:speedup_3way} shows the speedup of KV-stat
(\(n=3, lmc=16\), pipeline cycles) against all three baseline
variants.
The three columns tell three different stories:
the first is misleading, the second is real but driven by
score-matrix traffic, and the third is the honest comparison
after equal fusion.

\begin{table}[H]
\centering
\caption{Three-way prefill speedup: KV-stat \(n=3,\;lmc=16\)
cycles vs each baseline variant, wavefront model at \(512\,\)B/cycle.
Under the wavefront fill KV-stat is DRAM-bound at every \(T\), so it
matches the fused baseline (\(\sim\!1.01\times\)) throughout --- the
small-\(T\) deficit of the dwell model is gone.}
\label{tab:speedup_3way}
\begin{tabular}{rrrr}
\toprule
\(T\) & vs \(64{\times}64\) unfused & vs MAC-norm unfused & vs MAC-norm fused \\
\midrule
128   &   8.0$\times$ &  2.0$\times$ & \textbf{1.01$\times$} \\
512   &  31.9$\times$ &  5.0$\times$ & \textbf{1.01$\times$} \\
1,024 &  63.9$\times$ &  9.0$\times$ & \textbf{1.01$\times$} \\
2,048 & 127.8$\times$ & 17.0$\times$ & \textbf{1.01$\times$} \\
4,096 & 255.5$\times$ & 33.0$\times$ & \textbf{1.01$\times$} \\
8,192 & 511.0$\times$ & 64.9$\times$ & \textbf{1.01$\times$} \\
\bottomrule
\end{tabular}
\end{table}

Under the wavefront fill there is no crossover: KV-stat is memory-bound
at all \(T\) and, sharing the fused baseline's DRAM footprint exactly,
it tracks that baseline at \(\sim\!1.01\times\) everywhere. The large
multiples against the unfused baselines are unchanged --- the
\(\sim\!511\times\) figure is a MAC-count artifact of the compute-bound
\(64{\times}64\) array, and the \(\sim\!65\times\) is the
\(17.45\,\)GB score-matrix DRAM that the unfused baseline pays and both
fused designs eliminate. The honest equal-silicon comparison remains
\(\sim\!1\times\): KV-stat neither beats nor trails FlashAttention on
cycles when both eliminate the score matrix and run at the same
bandwidth.

\subsection{Prefill: KV-Stationary vs FlashAttention on Equal Silicon}
\label{sec:equalsilicon}

The primary comparison pairs KV-stat against a FlashAttention-style
fused baseline.
Both architectures eliminate score-matrix DRAM traffic and have
identical DRAM footprints; the comparison is therefore on
cycles and silicon area.
The FlashAttention reference uses a fixed \(64 \times 64\) array
(4{,}096 PEs, one MAC each); KV-stat scales with \(T\) as described in \S3.
All cycle counts below use the SCALE-Sim-validated wavefront fill
(\S\ref{sec:wavefront}) at a commodity bandwidth of \(512\,\)B/cycle.

The honest area metric is \textit{speedup per physical MAC}:
raw cycle speedup divided by the physical-MAC ratio. We deliberately
do \emph{not} normalize by PE \emph{count}, because each KV-stat PE packs
\(lmc\) lower MACs while each FlashAttention PE is a single MAC --- so
PE-count normalization understates KV-stat's silicon by a factor of \(lmc\).
We take \(\text{kv\_MACs} = \text{PE} \times lmc\) and
\(\text{flash\_MACs} = \text{PE} \times 1\). A per-MAC value above 1.0 means
KV-stat delivers more throughput per multiplier than FlashAttention.

\begin{table}[H]
\centering
\caption{Prefill --- KV-stat vs FlashAttention, both bandwidth regimes side by side
(wavefront model). \emph{Per-MAC} \(=\) raw speedup \(\div\) MAC ratio, where
MAC ratio \(= \tfrac{T}{64}\!\cdot\! lmc\) is KV-stat's physical multiplier count
vs the \(64{\times}64\) flash array. \textbf{DRAM-bound} is the commodity
\(512\,\)B/cycle point; \textbf{Compute-only} removes the DRAM floor (the
infinite-bandwidth / SCALE-Sim limit, reached with HBM-class bandwidth,
\S\ref{sec:wavefront}). The PE-\emph{count} ratio is \(T/64\) for all rows
(merge- and lmc-invariant), so the silicon difference lives entirely in MACs per
PE --- hence we normalize by MACs. \(Q\!\cdot\!K\) fill was SCALE-Sim-validated for
all rows except \(n=0, lmc=16\) at \(T \geq 4{,}096\) (operand-limit; analytical).
DRAM footprints are equal (both fused).}
\label{tab:prefill_results}
\begin{tabular}{lr rr rr r}
\toprule
& & \multicolumn{2}{c}{DRAM-bound (512\,B/cyc)} & \multicolumn{2}{c}{Compute-only} & \\
\cmidrule(lr){3-4}\cmidrule(lr){5-6}
Configuration & \(T\) & Raw & per-MAC & Raw & per-MAC & MAC ratio \\
\midrule
\(n=0,\;lmc=1\)  &  512 &   7.9$\times$ & 0.99$\times$ &   7.9$\times$ & 0.99$\times$ &     8$\times$ \\
\(n=0,\;lmc=1\)  & 2048 &  33.0$\times$ & 1.03$\times$ &  33.0$\times$ & 1.03$\times$ &    32$\times$ \\
\(n=0,\;lmc=1\)  & 8192 & 133.4$\times$ & 1.04$\times$ & 133.4$\times$ & 1.04$\times$ &   128$\times$ \\
\midrule
\(n=0,\;lmc=16\) &  512 &  33.0$\times$ & 0.26$\times$ &   124$\times$ & 0.97$\times$ &   128$\times$ \\
\(n=0,\;lmc=16\) & 2048 & 131.9$\times$ & 0.26$\times$ &   526$\times$ & 1.03$\times$ &   512$\times$ \\
\(n=0,\;lmc=16\) & 8192 & 527.4$\times$ & 0.26$\times$ & 2{,}133$\times$ & 1.04$\times$ & 2{,}048$\times$ \\
\midrule
\(n=3,\;lmc=16\) &  512 &  33.4$\times$ & 0.26$\times$ &   124$\times$ & 0.97$\times$ &   128$\times$ \\
\(n=3,\;lmc=16\) & 2048 & 133.7$\times$ & 0.26$\times$ &   775$\times$ & \textbf{1.51$\times$} &   512$\times$ \\
\(n=3,\;lmc=16\) & 8192 & 534.6$\times$ & 0.26$\times$ & 3{,}604$\times$ & \textbf{1.76$\times$} & 2{,}048$\times$ \\
\bottomrule
\end{tabular}
\end{table}

\begin{figure}[H]
    \centering
    \includegraphics[width=\textwidth]{fig20_prefill_both_regimes.png}
    \caption{Table~\ref{tab:prefill_results} visualized, both regimes side by side.
    \textbf{Left (a, c) --- DRAM-bound (512\,B/cyc):} raw speedup is large (up to
    \(\sim\!527\times\); \(n=0\) and \(n=3\) coincide), but per MAC the \(lmc = 16\)
    designs sit flat at \(0.26\times\) (below break-even, \(75\%\) idle), while lean
    \(lmc = 1\) hovers at \(\sim\!1.0\times\). \textbf{Right (b, d) --- compute-only
    (memory wall removed):} raw speedup jumps (to \(\sim\!2{,}133\times\) for \(n=0\),
    \(\sim\!3{,}604\times\) for \(n=3\)) and the curves diverge; per MAC the
    \(lmc = 16\) designs reach break-even (\(n=0 \approx 1.0\times\)) and merge
    extensions exceed it (\(n=3 \to 1.76\times\)).}
    \label{fig:prefill_permac}
\end{figure}

Three findings stand out from the DRAM-bound columns of
Table~\ref{tab:prefill_results} (left of Figure~\ref{fig:prefill_permac}), and they
revise the earlier (dwell-model) conclusions.

\textbf{(1) Large raw speedup, but it is bought with silicon.}
At \(512\,\)B/cycle the \(lmc = 16\) array finishes its compute well below the DRAM
floor, so its wall-clock speedup is large (up to \(527\times\) at \(T = 8{,}192\)).
Per physical MAC, however, it sits at only \(0.26\times\): it spends \(16\times\) the
multipliers of FlashAttention and, being memory-bound, leaves \(\sim\!75\%\) of them
idle (\S\ref{sec:codesign}). The lean \(lmc = 1\) array, though \(\sim\!4\times\)
slower in wall-clock, is compute-bound and lands at \(\sim\!1.0\times\) per MAC ---
more silicon-efficient. The per-MAC optimum is the balance point \(lmc \approx
\text{BW}/128\) (here \(lmc = 4\); Figure~\ref{fig:codesign}).

\textbf{(2) Merge extensions no longer help.}
Under the wavefront fill, \(n = 3, lmc = 16\) is indistinguishable from
\(n = 0, lmc = 16\) (both \(0.26\times\) per MAC, both memory-bound). Reshaping the
array to cut pipeline drain buys nothing once the design is DRAM-bound; merge
extensions only mattered in the dwell-model's compute-bound regime.

\textbf{(3) Per-MAC efficiency is sequence-length-invariant.}
Because prefill is DRAM-bound and both raw speedup and the MAC ratio scale linearly
with \(T\), per-MAC efficiency is flat across \(T\) (\(0.26\times\) for \(lmc = 16\),
\(\sim\!1.0\times\) for \(lmc = 1\)). The earlier ``advantage grows with \(T\)''
trend was an artifact of the dwell-model's drain term shrinking relative to compute.

\paragraph{The compute-bound limit (right-hand columns).}
The \emph{Compute-only} columns of Table~\ref{tab:prefill_results} --- and panels
(b, d) of Figure~\ref{fig:prefill_permac} --- remove the DRAM floor (the
infinite-bandwidth regime SCALE-Sim itself runs in, and the asymptote of the
bandwidth sweep in \S\ref{sec:wavefront}). Two reversals stand out. First,
\(lmc = 16\) is no longer wasteful: with every MAC fed, its per-MAC efficiency rises
from \(0.26\times\) to \(\sim\!1.0\times\), matching \(lmc = 1\) (per-MAC is
\(lmc\)-invariant once compute-bound). Second, merge extensions return ---
\(n = 3, lmc = 16\) reaches \(\mathbf{1.76\times}\) per MAC at \(T = 8{,}192\),
\emph{above} break-even --- because at equal MAC count its taller, narrower array
shortens the pipeline traversal (\(512 + 1024\) vs \(64 + 8192\)) and finishes its
compute sooner; the margin grows with \(T\) as the \(n = 0\) traversal it removes
grows. These are the conclusions a compute-bound (or HBM-provisioned) accelerator
would see; at commodity bandwidth the left-hand columns govern. The truth for a
given system is wherever its bandwidth places it on the \(lmc \approx \text{BW}/128\)
ladder (\S\ref{sec:wavefront}).

\subsection{Prefill: Comparison Against Unfused Baseline}

The full baseline characterization is presented in
\S\ref{subsec:baseline}.
In summary: the 65$\times$ cycle speedup against the MAC-normalized
unfused baseline at \(T = 8{,}192\) is entirely explained by the
17.45\,GB score-matrix DRAM that the unfused baseline incurs
(98.4\% of its total traffic) vs the 0.27\,GB footprint of KV-stat.
This advantage collapses to \(\sim\)1$\times$ when the baseline
adopts fused (FlashAttention-style) tiling, confirming that the
gain is not unique to the KV-stationary dataflow --- it is the
same score-matrix elimination that FlashAttention achieves in
software.

\subsection{Autoregressive Decode: Analytical Projection (Unvalidated)}

The analytical model projects a significant DRAM reduction for
autoregressive decode if the KV cache is held in on-chip SRAM
across generated tokens.
\textbf{This projection has not been validated with SCALE-Sim};
it is included to motivate the design direction described in
Section~\ref{sec:future}.

At \(T = 8{,}192\), \(H = 64\), \(d = 128\), FP16, the
per-token DRAM costs under the analytical model are:

\begin{center}
\begin{tabular}{lrl}
\toprule
Component & Per-step DRAM & Notes \\
\midrule
Conventional (KV reloaded) & 6.0\,MB & K + V + Q + scores + output \\
KV-stat model (KV in SRAM) & 32\,KB  & Q + output only \\
\midrule
Projected reduction & \textbf{188$\times$} & analytical only \\
\bottomrule
\end{tabular}
\end{center}

The 188$\times$ figure assumes the full KV cache fits in on-chip
SRAM throughout the generation run.
At \(T = 8{,}192\) this requires 128\,MB of on-chip SRAM ---
a substantial area commitment that motivates the two-pass
extension discussed in Section~\ref{sec:future}.

\subsection{Physical Cost}

The area savings implied by the speedup numbers require
substantial silicon.
For the best prefill configuration (\(n = 0, lmc = 16\)) at
\(T = 8{,}192\):

\begin{center}
\begin{tabular}{ll}
\toprule
Resource & Quantity \\
\midrule
Total PEs & 524{,}288 \\
Total MACs & 8{,}388{,}608 \\
On-chip K-buffer SRAM & 128\,MB \\
Estimated die area & \(\sim\)20\,mm\textsuperscript{2} \\
\bottomrule
\end{tabular}
\end{center}

The 128\,MB K-buffer alone is a significant on-chip memory
commitment --- roughly the L3 cache of a mid-range embedded
processor.
For the validated \(n = 0, lmc = 16, T = 2{,}048\) configuration,
requirements are 131{,}072 PEs and 34\,MB SRAM, which is
more tractable for edge deployment.

% ============================================================
\subsection{Corrected Pipeline Fill and Bandwidth Co-Design}
\label{sec:wavefront}

The timing model of \S\ref{sec:pipeline} originally charged the full
\(\text{column\_dwell} = 135\) cycles to every step of the pipeline fill. Because
each column has a dedicated lower MAC and a stationary \(K_i\), a single packet's
per-column dot products are independent and pipeline as a wavefront, so the
\(d\)-cycle reduction should be paid only \emph{once}. We tested which model is
physically correct using SCALE-Sim, which runs the \(Q\!\cdot\!K\) GEMM on a
cycle-accurate systolic array.

\paragraph{SCALE-Sim rules out the dwell model.}
We extend the Experiment~1 lower-MAC validation (\S\ref{sec:methodology}) from
\(lmc \in \{1,16\}\) to the full ladder \(lmc \in \{1,2,4,8,16\}\), to probe the
fill across the balance points of \S\ref{sec:codesign}. Across the resulting
\(40\) prefill \(Q\!\cdot\!K\) GEMMs (all fitting the \(500\,\)MB operand limit),
the measured per-column fill is \(2.1\)--\(3.0\) cycles for \(T \geq 512\) and is
independent of \(lmc\) (it is set by the \(H\times\text{eff\_cols}\) array
geometry). The dwell model's \(135\) cycles
per column overshoots by \(26\)--\(65\times\) and is decisively ruled out
(Figure~\ref{fig:fill_val}). The wavefront fill is the validated model; it is in fact
\emph{conservative} relative to SCALE-Sim. This validation covers the lower-MAC
compute only --- the same scope SCALE-Sim is used for throughout this work --- and runs
at infinite bandwidth to isolate compute.

\begin{figure}[H]
    \centering
    \includegraphics[width=0.8\textwidth]{fig13_scalesim_fill_validation.png}
    \caption{\textbf{SCALE-Sim validates the wavefront fill.} Measured per-column
    fill of the \(Q\!\cdot\!K\) stage clusters at \(\sim\!2.5\) cycles/column,
    versus the dwell model's \(135\). \(40\) cycle-accurate GEMM runs.}
    \label{fig:fill_val}
\end{figure}

\paragraph{Consequence: prefill is DRAM-bound.}
\label{sec:codesign}
With the wavefront fill, compute at \(T = 8{,}192\) drops far below the DRAM floor
of \(\sim\!525\)K cycles (moving \(273\,\)MB of \(Q\) and \(O\) at \(512\,\)B/cycle).
The array's compute scales as \(1/lmc\) while the DRAM floor is fixed by bandwidth,
so the two cross at a balance point \(lmc \approx \text{BW}/128\)
(Figure~\ref{fig:floor}). Below this point the design is memory-bound and additional
MACs sit idle --- at \(lmc = 16\) and \(512\,\)B/cycle the array is \(75\%\) idle.
The implication inverts the earlier recommendation: \(lmc\) must be sized to
bandwidth, not maximized. At commodity \(512\,\)B/cycle the efficient choice is
\(lmc \le 4\); \(lmc = 16\) is correct only with \(\sim\!2\,\)TB/s (HBM-class)
bandwidth, where it then runs \(4\times\) faster (Figures~\ref{fig:codesign},
\ref{fig:bwsweep}). Per physical MAC, every doubling of \(lmc\) past the balance
point halves efficiency, while at or below it efficiency is \(lmc\)-invariant and
\(lmc\) only buys lower wall-clock latency.

\begin{figure}[H]
    \centering
    \begin{subfigure}[t]{0.49\textwidth}
        \centering
        \includegraphics[width=\textwidth]{fig15_compute_vs_dram_floor.png}
        \caption{Compute (\(\propto 1/lmc\)) crosses below the fixed DRAM floor at
        \(lmc \approx \text{BW}/128\); past it the array idles.}
        \label{fig:floor}
    \end{subfigure}\hfill
    \begin{subfigure}[t]{0.49\textwidth}
        \centering
        \includegraphics[width=\textwidth]{fig14_lmc_bw_codesign.png}
        \caption{Per-MAC efficiency vs \(lmc\) per bandwidth. The balance point
        (\(\star\)) marches right as bandwidth grows.}
        \label{fig:codesign}
    \end{subfigure}
    \caption{\textbf{The \(lmc\)--bandwidth co-design.} \(T = 8{,}192\), single pass.}
\end{figure}

\begin{figure}[H]
    \centering
    \includegraphics[width=0.7\textwidth]{fig16_speedup_vs_bandwidth.png}
    \caption{\textbf{Speedup saturates at the compute floor.} For \(lmc = 16\),
    bandwidth helps up to \(\sim\!2\,\)TB/s (a \(4\times\) gain), then stops mattering
    --- past the knee the design is compute-bound and extra bandwidth is wasted.}
    \label{fig:bwsweep}
\end{figure}

\paragraph{The compute-bound limit.}
The asymptote of Figure~\ref{fig:bwsweep} --- the same configurations with the
memory wall removed (the compute-only, infinite-bandwidth regime SCALE-Sim itself
runs in) --- is reported in the \emph{Compute-only} columns of
Table~\ref{tab:prefill_results} (\S\ref{sec:equalsilicon}), side by side with the
DRAM-bound results. There, \(lmc = 16\) reaches break-even and merge extensions
exceed it; here we have shown \emph{why} --- that limit is only realized once
bandwidth feeds the array.

% ============================================================
\section{Discussion}

\subsection{What Makes the Advantage Real}

Both KV-stat and FlashAttention eliminate score-matrix DRAM traffic;
this is the shared insight underlying both architectures.
KV-stat's additional advantage is structural: it physically
prevents score materialization by keeping KV resident in the
compute fabric, rather than relying on tiling software to avoid it.
The analytical model suggests this property could extend to decode
--- where FlashAttention still reloads the full KV cache from DRAM
on every generated token --- but this has not been validated and
is left for future work.

\subsection{Design Space Findings}
\label{sec:findings}

The single most important design finding is that the right
configuration is \emph{bandwidth-dependent}, because under the
SCALE-Sim-validated wavefront fill (\S\ref{sec:pipeline}) prefill
compute is cheap and the design is normally DRAM-bound.

\textbf{Size \(lmc\) to bandwidth.}
Compute scales as \(1/lmc\) while the DRAM floor is fixed by bandwidth,
so the two balance at \(lmc \approx \text{BW}/128\) (\S\ref{sec:codesign}).
Below this point the array is memory-bound and the extra MACs sit idle
--- at \(lmc = 16\) on commodity \(512\,\)B/cycle DRAM, \(75\%\) of the
multipliers are stalled, and per-MAC efficiency is only \(0.26\times\).
At or above the balance point the design is compute-bound, per-MAC
efficiency is \(lmc\)-invariant (\(\sim\!1\times\)), and larger \(lmc\)
buys lower wall-clock latency. So \(lmc \le 4\) is correct at commodity
bandwidth; \(lmc = 16\) is justified only with HBM-class
(\(\sim\!2\,\)TB/s) bandwidth, where it then runs \(4\times\) faster.

\textbf{Merge extensions help only when compute-bound.}
Because merge extensions cut the column traversal (a \emph{compute}
term) without changing DRAM, they are inert once the design is
DRAM-bound: at \(512\,\)B/cycle \(n = 3\) and \(n = 0\) are
indistinguishable (both \(0.26\times\) per MAC). In the compute-bound
regime they return and are beneficial --- \(n = 3, lmc = 16\) reaches
\(1.76\times\) per MAC at \(T = 8{,}192\), above break-even
(Table~\ref{tab:prefill_results}, compute-only columns). This reverses
the earlier dwell-model claim that \(n = 0\) is universally
area-optimal.

\subsection{Validation Gaps}
\label{sec:gaps}

SCALE-Sim validation in this work covers the lower-MAC \(Q\!\cdot\!K\)
compute and its pipeline fill (\S\ref{sec:wavefront}); it runs at
infinite bandwidth to isolate compute and therefore does \emph{not}
validate the DRAM-bound cycle totals that govern the prefill results,
nor the \(lmc\)--bandwidth balance --- those rest on a roofline DRAM
model. Some large configurations (\(n = 0, lmc = 16\) at \(T \geq
4{,}096\)) exceed the \(500\,\)MB operand-matrix limit and so their fill
is extrapolated within the same validated model rather than directly
measured. Finally, the upper PE (exp lookup and \(V\) accumulation) and
the online-softmax state chain are modeled analytically throughout; a
hardware implementation may show different latency for the exp unit, and
the state-chain rate (assumed \(\le\) the measured \(\sim\)2.5\,cyc/col
fill) is not independently measured.

% ============================================================
\section{Conclusion}

This project built and validated an analytical performance model
of a KV-stationary 2D systolic array for MQA inference,
cross-checked against SCALE-Sim simulation across 29 configurations.

The principal findings are:

\begin{enumerate}[noitemsep]
    \item \textbf{Prefill (validated):} With the SCALE-Sim-validated
          wavefront fill, prefill is DRAM-bound and KV-stat matches a
          fused FlashAttention baseline on cycles at equal silicon
          (\(\sim\!1\times\)); per physical MAC it is below break-even
          for any over-provisioned \(lmc\) at commodity bandwidth.
          Both designs win their large multiples over an \emph{unfused}
          baseline (up to \(511\times\)) purely by eliminating the
          \(17\,\)GB score-matrix DRAM --- not from the dataflow itself.
          (An earlier dwell-fill model reported a \(\sim\)1.75--1.86$\times$
          per-area advantage; the wavefront correction removes it.)

    \item \textbf{Design recommendation (bandwidth-dependent):} Under the
          SCALE-Sim-validated wavefront fill (\S\ref{sec:wavefront}), prefill is
          DRAM-bound, and the lower-MAC count must be \emph{sized to bandwidth}
          rather than maximized: \(lmc \approx \text{BW}/128\) (so \(lmc \le 4\) at
          commodity \(512\,\)B/cycle, and \(lmc = 16\) only with HBM-class
          \(\sim\!2\,\)TB/s, where it then runs \(4\times\) faster). In this regime
          merge extensions (\(n = 3\)) collapse onto the flat array, since cutting
          drain buys nothing once memory-bound; their earlier advantage holds only
          where the design is compute-bound (small \(lmc\) or high bandwidth).

    \item \textbf{Decode (future work):} Autoregressive decode was not
          validated in this work.
          The analytical model projects a 188$\times$ per-token DRAM
          reduction if the KV cache is held on-chip, but this requires
          up to 128\,MB of SRAM.
          A two-pass decode scheme is the proposed path to making
          this practical at lower chip area (Section~\ref{sec:future}).
\end{enumerate}

\textbf{The central limitation of this architecture is physical scaling.}
The headline speedup numbers --- 165$\times$ at \(T = 8{,}192\) ---
require \textbf{524{,}288 PEs} and \textbf{128\,MB of on-chip SRAM}
just for the K-buffer, before routing, control logic, or the upper
Running Attention PEs are counted.
At the validated \(T = 2{,}048\) operating point the requirements
are more modest (131K PEs, 34\,MB SRAM), but even here the array
is 32$\times$ larger in PEs than the FlashAttention reference while
matching it on cycles. The wavefront analysis sharpens this point:
because prefill is DRAM-bound, the large raw speedups buy no
per-silicon advantage --- per physical MAC the architecture sits at
break-even (\(\sim\!1\times\) for a bandwidth-matched \(lmc\)) and
strictly below it for any over-provisioned \(lmc\). No configuration
in the tested space exceeds \(\sim\!1\times\) per MAC, and the PE count
and SRAM still double with every doubling of~\(T\).

In short: the KV-stationary dataflow is architecturally sound and its
score-matrix elimination is real, but on equal silicon and equal
bandwidth it matches rather than beats FlashAttention, and the chip
cost of supporting the raw speedup numbers grows too fast with
sequence length to be practical at large~\(T\) without the kind of
multi-pass scheduling described in Section~\ref{sec:future}.

% ============================================================
\section{Future Work: Two-Pass Decode}
\label{sec:future}

The primary barrier to practical decode deployment is on-chip SRAM.
At \(T = 8{,}192\) the full KV cache requires 128\,MB of on-chip
storage --- prohibitive for edge or embedded targets.
A \textbf{two-pass decode scheme} addresses this directly.

\textbf{Concept.}
Rather than holding the entire KV cache on-chip across all tokens,
the sequence is partitioned into \(P\) equal slices of
\(T / P\) tokens each.
Each decode step makes \(P\) passes over the array, one per slice,
accumulating the running softmax state \((m, l, O)\) across passes
using the same online-softmax merge formula used by the merge-tree
for prefill.
After all \(P\) passes the final output is produced identically
to the single-pass case.

Figure~\ref{fig:reuse_model} shows the \(P = 2\) case concretely.
In Pass~1 the query stream enters the first slice of KV columns
(\texttt{KV0--KV15}) and produces a partial softmax state
\texttt{O\_partial} --- the running \((m, l, O)\) vector after
attending over the first \(T/2\) tokens.
This state is saved to a small output register (green box).
In Pass~2 the same query stream re-enters, this time through
the second slice (\texttt{KV16--KV31}).
Before the stream reaches the first PE of the second slice,
the saved \texttt{O\_partial} state is re-injected via
\texttt{PE0} --- a single merge PE that combines the partial
state with the incoming running state using the standard
online-softmax update.
The second slice then continues accumulating from the correct
starting point, and the final column produces
\texttt{O\_final}: the correctly normalised attention output
over all \(T\) tokens.

The hardware change relative to a single-pass design is minimal:
one additional merge PE at the slice boundary and a small
register to hold the \((m, l, O)\) state between passes
(\(d + 2\) elements = 258 bytes at FP16).
The array itself is unchanged --- the same \(H \times T/P\)
physical columns are reused across both passes.

\begin{figure}[H]
    \centering
    \includegraphics[width=0.82\textwidth]{ReuseModel.png}
    \caption{\textbf{Silicon Reduction and Re-Use Model (\(P = 2\) passes).}
    Pass~1: the query stream attends over the first KV slice (\texttt{KV0--KV15}),
    producing a partial softmax state \texttt{O\_partial} that is saved.
    Pass~2: the same query re-enters and attends over the second slice
    (\texttt{KV16--KV31}); the saved partial state is re-injected via \texttt{PE0}
    so the online softmax continues from the correct running maximum and denominator.
    \texttt{O\_final} is the correctly normalised output over all \(T\) tokens.
    The array uses half the columns (and half the SRAM) of a single-pass design
    at the cost of running each decode token twice.}
    \label{fig:reuse_model}
\end{figure}

\textbf{Silicon reduction and the latency trade-off.}
The re-use model applies to both prefill and decode.
In either phase, using \(P\) passes over a \(T/P\)-column
array instead of one pass over a \(T\)-column array reduces
PE count and SRAM by \(P\times\):

\[
\text{columns} = T / P, \qquad
\text{SRAM\_required} = \frac{2 \cdot T \cdot d \cdot \text{bpe}}{P}
\]

The cost is \textbf{latency}, not DRAM.
Each pass requires the query stream to traverse the array
a second time, and at each pass boundary the partial softmax
state \((m, l, O)\) must be written to a register and
re-injected via \texttt{PE0} at the start of the next pass.
This save-and-reinject overhead is small --- 258 bytes at
FP16 --- but the \(P\times\) pipeline repetition is a real
latency multiplier for both prefill and decode.
There is no DRAM penalty: the KV slices are already resident
on-chip and no additional DRAM traffic is incurred by
the re-use itself.

\begin{center}
\begin{tabular}{lrrll}
\toprule
Passes \(P\) & Columns & On-chip SRAM & Cycle cost & Applies to \\
\midrule
1 & \(T\)     & 128\,MB & \(1\times\) & prefill + decode \\
2 & \(T/2\)   &  64\,MB & \(2\times\) + reinject & prefill + decode \\
4 & \(T/4\)   &  32\,MB & \(4\times\) + reinject & prefill + decode \\
8 & \(T/8\)   &  16\,MB & \(8\times\) + reinject & prefill + decode \\
\bottomrule
\end{tabular}
\end{center}

\textbf{Why this is worth building.}
The primary value is silicon area reduction, not DRAM savings.
The fundamental scaling problem identified in \S7 is that
PE count and SRAM grow linearly with \(T\): at \(T = 8{,}192\)
the single-pass design requires 524K PEs and 128\,MB SRAM.
With \(P = 4\) passes the same computation fits in a
\(T/4 = 2{,}048\)-column array --- 131K PEs and 32\,MB SRAM
--- which is the validated operating point from the prefill
experiments, requiring only \(4\times\) more cycles per token.
In effect, the re-use model lets the designer trade latency
for chip area at a known, controllable ratio, making the
architecture viable at sequence lengths that would otherwise
be physically impractical.

\textbf{Open questions for future validation.}
\begin{itemize}[noitemsep]
    \item SCALE-Sim cycle validation of a single-pass decode step
          (currently divergent due to pipeline drain, as noted in \S\ref{sec:methodology}).
    \item Extending the analytical pipeline model to multi-pass decode,
          including the inter-pass softmax merge cost.
    \item Determining the optimal \(P\) that balances SRAM budget
          against acceptable decode latency for the target workload.
\end{itemize}

% ============================================================
\section{References}

\printbibliography[heading=none]

\end{document}
